%
%****************************************************************************************************
%========================= B E G I N   D O C U M E N T   A D D E N D U M =====================
%****************************************************************************************************
%
\section{Anhang}
\label{sec:anhang}

\subsection{Herleitung der Speicherdauerformel}

Die in Abschnitt~\ref{sec:denormalisierung} verwendete Formel zur Berechnung der Speicherauffüllungsdauer wurde wie folgt hergeleitet:

Die durchschnittliche Größe eines Datensatzes ($s$) ergibt sich aus:

\begin{equation}
    s = \frac{G_t}{E_t}
\end{equation}

wobei $G_t$ die Größe eines Backups und $E_t$ die Gesamtanzahl der Erkennungen zu diesem Zeitpunkt $t$ bezeichnen.
Die tägliche Erkennungsrate $r$ berechnet sich folgendermaßen:

\begin{equation}
    r = \frac{E_t}{D_t}
\end{equation}

wobei $D_t$ die Anzahl der bisherigen Aufzeichnungstage darstellt.
Daraus lässt sich die Anzahl der Tage $T$ bestimmen, bis eine gegebene Speicherkapazität $S_{\max}$ erreicht ist:

\begin{equation}
    T = \frac{S_{\max}}{s \cdot r}
\end{equation}

Ausgehend von den oben aufgelisteten Formeln ergibt sich durch Einsetzen von $s$ und $r$ in $T$:

\begin{equation}
    T = \frac{S_{\max}}{\left(\frac{G_t}{E_t}\right)\left(\frac{E_t}{D_t}\right)}
\end{equation}

Nach Kürzen von $E_t$ folgt:

\begin{equation}
    T = \frac{S_{\max}}{\frac{G_t}{D_t}}
\end{equation}

wobei $\frac{G_t}{D_t}$ die durchschnittliche Speicherzunahme pro Tag beschreibt.
Damit ergibt sich die vereinfachte Formel:

\begin{equation}
    T = \frac{S_{\max} \cdot D_t}{G_t}
\end{equation}

\newpage
% 1. Force the layout to start on an EVEN page (the left side of a spread)
\ifodd\value{page}\hbox{}\newpage\fi
\subsection{Sequenzdiagramm des Kennzeichenerkennungsprozess}

% --- LEFT PAGE (P1) ---
\begin{tikzpicture}[remember picture, overlay]
% Anchor to the spine, shifted down by 2cm to avoid covering the subsection heading
\node[anchor=east, inner sep=0pt, yshift=-2cm, xshift=-22mm] at (current page.east) {
\adjustbox{
        trim={0pt} {0.44\height} {0pt} {0pt},
        clip,
        angle=90               
    }{%
% Scaled to leave margin space for the heading and caption
\includegraphics[width=0.78\paperheight, height=2\paperwidth, keepaspectratio]{data_collection.png}%
    }
};
\end{tikzpicture}

% Pushes the caption to the bottom of the left page
\vfill
\captionof{figure}{Data Collection Flow}
\label{fig:data_collection_flow}
\vspace{1cm}

\newpage

% --- RIGHT PAGE (P2) ---
\begin{tikzpicture}[remember picture, overlay]
% Anchor to the spine, shifted down by the exact same amount to ensure alignment
\node[anchor=west, inner sep=0pt, yshift=-2cm, xshift=30mm] at (current page.west) {
\adjustbox{
            trim={0pt} {0pt} {0pt} {0.56\height}, 
            clip,
            angle=90               
        }{%
\includegraphics[width=0.78\paperheight, height=2\paperwidth, keepaspectratio]{data_collection.png}%
        }
    };
\end{tikzpicture}

\newpage
\subsection{Vollständige Verarbeitungsmethode (process\_vehicle\_detection)}

Die folgende Methode zeigt den vollständigen Ablauf der Hintergrundverarbeitung einer Fahrzeugerkennung im Data Collection Service:

\begin{lstlisting}[language=Python, caption={Vollständige Verarbeitungsmethode}, label={lst:process_full}]
def process_vehicle_detection(camera_name: str, detection_time: datetime):
    """background task for vehicle detection handling

    Args:
        camera_name (str): name of the camera that called the webhook
        detection_time (datetime): time the event was triggered
    """
    try:
        with get_db() as db:
            # Authenticate client
            sid = camera_service.authenticate_client(
                host=settings.synology_host,
                username=settings.synology_username,
                password=settings.synology_password,
            )

            # Get camera data
            cameras = camera_service.get_camera_data(
                host=settings.synology_host, sid=sid)

            # Find target camera
            target_camera = camera_service.get_camera_by_name(
                cameras=cameras, camera_name=camera_name)

            # Get camera snapshot
            frame = camera_service.get_camera_snapshot(
                host=settings.synology_host, sid=sid, 
                camera=target_camera)

            image_data = frame.content
            filename = f'observation_{detection_time}.jpg'
            filepath = os.path.join('/app/snapshots', filename)

            # Save image is enabled
            if settings.save_images_for_debug:
                try:
                    with open(filepath, 'wb') as f:
                        f.write(image_data)
                    logger.info(f'Snapshot saved to {filepath}')
                except Exception as e:
                    logger.exception(f'Failed to save image: {e}')

            # Send image to api
            result = plate_service.send_to_api(
                api_key=settings.api_key,
                image_data=image_data,
                camera_name=camera_name,
                service_url=settings.plate_recognizer_service_url,
            )
            if not result:
                logger.info(
                    'Plate Recognizer returned no actual observations')
                return
            logger.debug(f'Plate Recognizer results: {result}')

            # Add result to db
            observations_to_create = db_handler.new_observation(
                reader_result=result, 
                detection_timestamp=detection_time
            )

            for observation_data in observations_to_create:
                observation_data = \
                    country_handler.get_municipality_and_fix_country(
                        observation=observation_data)
                observation_data = db_handler.hash_plate(
                    observation=observation_data)

                if not db_handler.check_for_duplicates(
                    db=db,
                    observation=observation_data,
                    current_detection_time=detection_time,
                    interval=timedelta(
                        seconds=settings.interval_seconds),
                ):
                    db_handler.create_observation_entry(
                        db=db, observation=observation_data)

    except (CameraException, PlateRecognizerException, 
            DatabaseException) as e:
        logger.exception(f'{type(e).__name__}: {e}')
    except Exception as e:
        logger.exception(
            f'Unexpected error during background processing: {e}')
    return
\end{lstlisting}

\newpage
\subsection{Datenbank-Initialisierungsskript (init.sql)}

Das folgende SQL-Skript wird beim ersten Start des Systems durch den DB-Prestart-Container ausgeführt und erstellt die Schemas, Benutzer und Berechtigungen:

\begin{lstlisting}[language=SQL, caption={Datenbank-Initialisierungsskript (init.sql)}, label={lst:init_sql}]
-- Create users with environment variable passwords
CREATE USER notification_user 
    WITH PASSWORD :NOTIFICATION_DB_PASSWORD;
CREATE USER data_collection_user 
    WITH PASSWORD :DATA_COLLECTION_DB_PASSWORD;
CREATE USER analytics_user 
    WITH PASSWORD :ANALYTICS_DB_PASSWORD;

-- Grant database connection privileges
GRANT CONNECT ON DATABASE :DB_NAME TO data_collection_user;
GRANT CONNECT ON DATABASE :DB_NAME TO analytics_user;
GRANT CONNECT ON DATABASE :DB_NAME TO notification_user;

-- Create schemas
CREATE SCHEMA IF NOT EXISTS :DATA_COLLECTION_SCHEMA;
CREATE SCHEMA IF NOT EXISTS :ANALYTICS_SCHEMA;
CREATE SCHEMA IF NOT EXISTS :NOTIFICATION_SCHEMA;

-- Data Collection Service permissions 
-- (writes to :DATA_COLLECTION_SCHEMA)
GRANT ALL ON SCHEMA :DATA_COLLECTION_SCHEMA 
    TO data_collection_user;
GRANT ALL PRIVILEGES ON ALL TABLES 
    IN SCHEMA :DATA_COLLECTION_SCHEMA TO data_collection_user;
GRANT ALL PRIVILEGES ON ALL SEQUENCES 
    IN SCHEMA :DATA_COLLECTION_SCHEMA TO data_collection_user;
ALTER DEFAULT PRIVILEGES IN SCHEMA :DATA_COLLECTION_SCHEMA 
    GRANT ALL ON TABLES TO data_collection_user;
ALTER DEFAULT PRIVILEGES IN SCHEMA :DATA_COLLECTION_SCHEMA 
    GRANT ALL ON SEQUENCES TO data_collection_user;

-- Analytics Service permissions 
-- (reads from :DATA_COLLECTION_SCHEMA, 
--  writes to :ANALYTICS_SCHEMA)
GRANT USAGE ON SCHEMA :DATA_COLLECTION_SCHEMA 
    TO analytics_user;
GRANT SELECT ON ALL TABLES 
    IN SCHEMA :DATA_COLLECTION_SCHEMA TO analytics_user;
ALTER DEFAULT PRIVILEGES IN SCHEMA :DATA_COLLECTION_SCHEMA 
    GRANT SELECT ON TABLES TO analytics_user;

GRANT ALL ON SCHEMA :ANALYTICS_SCHEMA TO analytics_user;
GRANT ALL PRIVILEGES ON ALL TABLES 
    IN SCHEMA :ANALYTICS_SCHEMA TO analytics_user;
GRANT ALL PRIVILEGES ON ALL SEQUENCES 
    IN SCHEMA :ANALYTICS_SCHEMA TO analytics_user;
ALTER DEFAULT PRIVILEGES IN SCHEMA :ANALYTICS_SCHEMA 
    GRANT ALL ON TABLES TO analytics_user;
ALTER DEFAULT PRIVILEGES IN SCHEMA :ANALYTICS_SCHEMA 
    GRANT ALL ON SEQUENCES TO analytics_user;

-- Notification Service permissions 
-- (writes to :NOTIFICATION_SCHEMA)
GRANT ALL ON SCHEMA :NOTIFICATION_SCHEMA 
    TO notification_user;
GRANT ALL PRIVILEGES ON ALL TABLES 
    IN SCHEMA :NOTIFICATION_SCHEMA TO notification_user;
GRANT ALL PRIVILEGES ON ALL SEQUENCES 
    IN SCHEMA :NOTIFICATION_SCHEMA TO notification_user;
ALTER DEFAULT PRIVILEGES IN SCHEMA :NOTIFICATION_SCHEMA 
    GRANT ALL ON TABLES TO notification_user;
ALTER DEFAULT PRIVILEGES IN SCHEMA :NOTIFICATION_SCHEMA 
    GRANT ALL ON SEQUENCES TO notification_user;

-- Set search paths for each user
ALTER ROLE data_collection_user 
    SET search_path = :DATA_COLLECTION_SCHEMA,public;
ALTER ROLE analytics_user 
    SET search_path = :ANALYTICS_SCHEMA,
                      :DATA_COLLECTION_SCHEMA,public;
ALTER ROLE notification_user 
    SET search_path = :NOTIFICATION_SCHEMA,public;
\end{lstlisting}

\newpage
\subsection{Stundenmitschrift}

Die folgende Tabelle dokumentiert den zeitlichen Aufwand für diese Diplomarbeit.

\begin{longtable}{|c|l|r|}
    \hline
    \textbf{Datum}  & \textbf{Tätigkeit}                             & \textbf{Zeitaufwand in h} \\
    \hline
    \endfirsthead
    \hline
    \textbf{Datum}  & \textbf{Tätigkeit}                             & \textbf{Zeitaufwand in h} \\
    \hline
    \endhead
    \hline
    \endfoot
    17.06.2025      & Arbeit am Pflichtenheft                        & 2                         \\
    24.06.2025      & Arbeit am Pflichtenheft                        & 2                         \\
    30.06.2025      & Aufbau Testrig \& Pflichtenheft                & 6                         \\
    01.07.2025      & Pflichtenheft \& Kamera                        & 7                         \\
    04.07.2025      & Prototyping                                    & 6                         \\
    07.07.2025      & Schnittstelle zu Synology                      & 6                         \\
    08.07.2025      & Setup third party API + LDAP Hilfe             & 7                         \\
    09.07.2025      & LDAP + API Anbindung Erweiterung               & 5                         \\
    10.07.2025      & Verbesserung Kennzeichen \& Landerkennung      & 7                         \\
    11.07.2025      & Kubernetes Test                                & 5                         \\
    14.07.2025      & Repo Setup                                     & 5                         \\
    15.07.2025      & Architektur Ausarbeitung \& Repo Setup         & 4                         \\
    16.07.2025      & Architektur \& Repo                            & 6                         \\
    18.07.2025      & Portainer Einlesen \& Repo finalisieren        & 5                         \\
    21.07.2025      & Setup finalisieren                             & 7                         \\
    22.07.2025      & Alembic                                        & 5                         \\
    23.07.2025      & Alembic Migration fixed                        & 2                         \\
    24.07.2025      & Aufbau Vorbereitung                            & 2                         \\
    25.07.2025      & Implementierung Datenspeicherung               & 7                         \\
    28.07.2025      & Alembic Verbesserungen \& Strukturänderung     & 6                         \\
    29.07.2025      & Alembic Verbesserte Integrierung mit k8s       & 5                         \\
    30.07.2025      & License plate fix                              & 5                         \\
    31.07.2025      & Aufbau und Inbetriebnahme Kamera               & 6                         \\
    01.08.2025      & Testkonzept + Tests für Data Collection        & 5                         \\
    04.08.2025      & Start Notification Service + VPN Zugriff Setup & 7                         \\
    05.08.2025      & Notification Service Arbeit                    & 5                         \\
    06.08.2025      & Notification Service Api                       & 5                         \\
    07.08.2025      & Integrierung der Plate Recognizer SDK          & 5                         \\
    08.08.2025      & NAS-Integration des Vorläufigen Systems        & 4                         \\
    29.10.2025      & Code-Stil Fixes, Build Improvements            & 1                         \\
    30.10.2025      & DevOps (Compose + Github Actions)              & 2                         \\
    01.11.2025      & MkDocs Setup + Deploy                          & 3                         \\
    07.11.2025      & Health Calls, Tests, Fixes                     & 2                         \\
    09.11.2025      & CI/CD Improvements + User Management Tests     & 3                         \\
    20.11.2025      & Grafana Implementiert                          & 3                         \\
    26.01.2026      & Notifications Service fertig                   & 4                         \\
    26.01.2026      & Dokumentation Struktur erstellt                & 1                         \\
    01.02.2026      & Pre-Content fertig                             & 2                         \\
    03.02.2026      & Einleitung fertig geschrieben                  & 3                         \\
    05.02.2026      & Theoretische Grundlagen start                  & 2                         \\
    06.02.2026      & Theoretische Grundlagen fertig                 & 2                         \\
    07.02.2026      & Hardwareauswahl fertig                         & 4                         \\
    08.02.2026      & Hardwareauswahl fixes + Start Systementwurf    & 3                         \\
    11.02.2026      & Systementwurf und Architektur geschrieben      & 4                         \\
    12.02.2026      & Formel für Speicherberechnung hinzugefügt      & 1                         \\
    13.02.2026      & Kapitel Implementierung hinzugefügt            & 6                         \\
    18.02.2026      & Kapitel Infrastruktur geschrieben              & 7                         \\
    19.02.2026      & Qualitätssicherung und Testkapitel fertig      & 5                         \\
    20.02.2026      & Fazit und Ausblick fertig + Korrekturen        & 3                         \\
    22.02.2026      & LaTeX Portierung                               & 5                         \\
    \hline
    \textbf{Gesamt} &                                                & \textbf{215}              \\
    \hline
\end{longtable}


%
\newpage 
\listoffigures 
%
\glsaddallunused
\newpage
\printnoidxglossary[type=\acronymtype, title={Abkürzungsverzeichnis}, sort=letter]
%
\newpage
%
% Literaturverzeichnis
%
\printbibliography[title={Literaturverzeichnis}]

% \begin{thebibliography}{99}
%
%     % --- Kapitel 1: Einleitung ---
%     \bibitem{steiermark_ausflugsziele}
%     STEIERMARK TOURISMUS, Die meistbesuchten Ausflugsziele der Steiermark. In: \url{https://www.steiermark.com/de/Magazin/Die-meistbesuchten-Ausflugsziele-der-Steiermark_mad_45932860} (am 23.2.2026).
%
%     \bibitem{ghg_scope3}
%     GREENHOUSE GAS PROTOCOL, Corporate Value Chain (Scope 3) Standard. In: \url{https://ghgprotocol.org/corporate-value-chain-scope-3-standard} (am 23.2.2026).
%
%     % --- Kapitel 2: Theoretische Grundlagen ---
%     \bibitem{docker_overview}
%     DOCKER, Docker overview. In: \url{https://docs.docker.com/get-started/docker-overview/} (am 23.2.2026).
%
%     \bibitem{docker_dockerfile}
%     DOCKER, Dockerfile reference. In: \url{https://docs.docker.com/build/concepts/dockerfile/} (am 23.2.2026).
%
%     \bibitem{docker_cache}
%     DOCKER, Docker build cache. In: \url{https://docs.docker.com/build/cache/} (am 23.2.2026).
%
%     \bibitem{docker_compose}
%     DOCKER, Docker Compose application model. In: \url{https://docs.docker.com/compose/intro/compose-application-model/} (am 23.2.2026).
%
%     \bibitem{docker_hub}
%     DOCKER, Docker Hub. In: \url{https://www.docker.com/products/docker-hub/} (am 23.2.2026).
%
%     \bibitem{aws_microservices}
%     AMAZON WEB SERVICES, What are Microservices? In: \url{https://aws.amazon.com/microservices/} (am 23.2.2026).
%
%     \bibitem{oso_microservices}
%     OSO, Microservices Best Practices. In: \url{https://www.osohq.com/learn/microservices-best-practices} (am 23.2.2026).
%
%     \bibitem{platerecognizer_alpr}
%     PLATE RECOGNIZER, What is ALPR?. In: \url{https://platerecognizer.com/what-is-alpr/} (am 23.2.2026).
%
%     \bibitem{senstar_alpr}
%     SENSTAR, Wie ALPR funktioniert. In: \url{https://senstar.com/de/senstarpedia/wie-alpr-funktioniert/} (am 23.2.2026).
%
%     \bibitem{platerecognizer_mmc}
%     PLATE RECOGNIZER, Vehicle Make Model Recognition with Color. In: \url{https://platerecognizer.com/vehicle-make-model-recognition-with-color/} (am 23.2.2026).
%
%     \bibitem{platerecognizer_snapshot}
%     PLATE RECOGNIZER, Snapshot SDK. In: \url{https://platerecognizer.com/snapshot/} (am 23.2.2026).
%
%     % --- Kapitel 3: Hardware & Technologie ---
%     \bibitem{baseus_ip_ratings}
%     BASEUS, A Guide to IP65, IP66, and IP67 Ratings for Outdoor Security Cameras. In: \url{https://security.baseus.com/en-eu/blogs/content/a-guide-to-ip65-ip66-and-ip67-ratings-for-outdoor-security-cameras} (am 23.2.2026).
%
%     \bibitem{wiki_ip_camera}
%     WIKIPEDIA, IP camera. In: \url{https://en.wikipedia.org/wiki/IP_camera} (am 23.2.2026).
%
%     \bibitem{synology_bc500}
%     SYNOLOGY, BC500. In: \url{https://www.synology.com/en-eu/products/BC500} (am 23.2.2026).
%
%     \bibitem{platerecognizer_camera_setup}
%     PLATE RECOGNIZER, Camera Setup for Best ANPR. In: \url{https://platerecognizer.com/camera-setup-for-best-anpr/} (am 23.2.2026).
%
%     \bibitem{aws_nas}
%     AMAZON WEB SERVICES, What is NAS? In: \url{https://aws.amazon.com/what-is/nas/} (am 23.2.2026).
%
%     \bibitem{synology_surveillance}
%     SYNOLOGY, Surveillance Station. In: \url{https://www.synology.com/en-global/surveillance} (am 23.2.2026).
%
%     \bibitem{synology_dva}
%     SYNOLOGY, Deep Video Analytics (DVA). In: \url{https://www.synology.com/en-global/products/DVA} (am 23.2.2026).
%
%     \bibitem{synology_lpr}
%     SYNOLOGY, License Plate Recognition (LPR). In: \url{https://kb.synology.com/de-de/UG/DVA_License_Plate_Recognition_AG/7} (am 23.2.2026).
%
%     \bibitem{platerecognizer_countries}
%     PLATE RECOGNIZER, Supported Countries. In: \url{https://platerecognizer.com/countries/} (am 23.2.2026).
%
%     \bibitem{platerecognizer_results}
%     PLATE RECOGNIZER, ALPR Results. In: \url{https://platerecognizer.com/alpr-results} (am 23.2.2026).
%
%     \bibitem{fastapi}
%     RAMÍREZ Sebastián, FastAPI. In: \url{https://fastapi.tiangolo.com/} (am 23.2.2026).
%
%     \bibitem{jetbrains_frameworks}
%     JETBRAINS, Django, Flask, FastAPI -- Which Python Web Framework to Choose?. In: \url{https://blog.jetbrains.com/pycharm/2025/02/django-flask-fastapi/} (am 23.2.2026).
%
%     \bibitem{postgresql}
%     THE POSTGRESQL GLOBAL DEVELOPMENT GROUP, PostgreSQL -- The World's Most Advanced Open Source Relational Database. In: \url{https://www.postgresql.org/} (am 23.2.2026).
%
%     \bibitem{alembic}
%     BAYER Michael, Alembic -- A database migration tool for SQLAlchemy. In: \url{https://alembic.sqlalchemy.org/en/latest/} (am 23.2.2026).
%
%     % --- Kapitel 4: Systementwurf & Architektur ---
%     \bibitem{twelve_factor}
%     WIGGINS Adam, The Twelve-Factor App. In: \url{https://12factor.net/} (am 23.2.2026).
%
%     \bibitem{nextjs}
%     VERCEL, Next.js. In: \url{https://nextjs.org/} (am 23.2.2026).
%
%     \bibitem{postgresql_schemas}
%     THE POSTGRESQL GLOBAL DEVELOPMENT GROUP, PostgreSQL Schemas. In: \url{https://www.postgresql.org/docs/current/ddl-schemas.html} (am 23.2.2026).
%
%     \bibitem{so_backup_size}
%     STACK OVERFLOW, Is database backup size the same as database size?. In: \url{https://stackoverflow.com/questions/35480650/is-database-back-up-size-is-same-as-database-size} (am 23.2.2026).
%
%     \bibitem{sqlskills_backup}
%     SQLSKILLS, Trending Database Growth from Backups. In: \url{https://www.sqlskills.com/blogs/erin/trending-database-growth-from-backups/} (am 23.2.2026).
%
%     \bibitem{wiki_polp}
%     WIKIPEDIA, Principle of least privilege. In: \url{https://en.wikipedia.org/wiki/Principle_of_least_privilege} (am 23.2.2026).
%
%     \bibitem{dsgvo_personenbezogen}
%     DSGVO-GESETZ.DE, Personenbezogene Daten. In: \url{https://dsgvo-gesetz.de/themen/personenbezogene-daten/} (am 23.2.2026).
%
%     % --- Kapitel 5: Implementierung ---
%     \bibitem{synology_webhooks}
%     SYNOLOGY, How do I send Surveillance Station data to third-party systems via webhooks?. In: \url{https://kb.synology.com/en-global/Surveillance/tutorial/How_do_I_send_SS_data_to_third-party_systems_via_webhooks} (am 23.2.2026).
%
%     \bibitem{synology_api}
%     SYNOLOGY, Surveillance Station Web API. In: \url{https://global.download.synology.com/download/Document/Software/DeveloperGuide/Package/SurveillanceStation/All/enu/Surveillance_Station_Web_API.pdf} (am 23.2.2026).
%
%     \bibitem{pydantic_settings}
%     PYDANTIC, Settings Management. In: \url{https://docs.pydantic.dev/latest/concepts/pydantic_settings/#installation} (am 23.2.2026).
%
%     \bibitem{tenacity}
%     MERGIFY, tenacity -- Retrying library for Python. In: \url{https://tenacity.readthedocs.io/en/latest/} (am 23.2.2026).
%
%     \bibitem{platerecognizer_api}
%     PLATE RECOGNIZER, Snapshot API Reference. In: \url{https://guides.platerecognizer.com/docs/snapshot/api-reference} (am 23.2.2026).
%
%     \bibitem{oesterreich_kfz}
%     OESTERREICH.GV.AT, Kennzeichnung von Kraftfahrzeugen. In: \url{https://www.oesterreich.gv.at/de/themen/mobilitaet/kfz/5/1} (am 23.2.2026).
%
%     \bibitem{wiki_slovenia_plates}
%     WIKIPEDIA, Vehicle registration plates of Slovenia. In: \url{https://en.wikipedia.org/wiki/Vehicle_registration_plates_of_Slovenia} (am 23.2.2026).
%
%     \bibitem{mailtrap_python_email}
%     MAILTRAP, How to Send an Email Using Python (Gmail). In: \url{https://mailtrap.io/blog/python-send-email-gmail/} (am 23.2.2026).
%
%     \bibitem{mailjet_smtp_relay}
%     MAILJET, SMTP Relay -- What it is and How it Works. In: \url{https://www.mailjet.com/blog/email-best-practices/what-is-an-smtp-relay/} (am 23.2.2026).
%
%     \bibitem{grafana}
%     GRAFANA LABS, Grafana -- The open observability platform. In: \url{https://grafana.com/} (am 23.2.2026).
%
%     \bibitem{grafana_provisioning}
%     GRAFANA LABS, Provision Grafana. In: \url{https://grafana.com/docs/grafana/latest/administration/provisioning/} (am 23.2.2026).
%
%     \bibitem{grafana_configure}
%     GRAFANA LABS, Configure Grafana. In: \url{https://grafana.com/docs/grafana/latest/setup-grafana/configure-grafana/} (am 23.2.2026).
%
%     % --- Kapitel 6: Infrastruktur ---
%     \bibitem{docker_buildx}
%     DOCKER, Docker Buildx. In: \url{https://docs.docker.com/buildx/working-with-buildx/} (am 23.2.2026).
%
%     \bibitem{portainer}
%     PORTAINER, Container Management Software. In: \url{https://docs.portainer.io/} (am 23.2.2026).
%
%     \bibitem{redhat_cicd}
%     RED HAT, What is CI/CD?. In: \url{https://www.redhat.com/en/topics/devops/what-is-ci-cd} (am 23.2.2026).
%
%     \bibitem{github_actions}
%     GITHUB, GitHub Actions Documentation. In: \url{https://docs.github.com/en/actions} (am 23.2.2026).
%
%     \bibitem{ruff_formatter}
%     ASTRAL, Ruff Formatter. In: \url{https://docs.astral.sh/ruff/formatter/} (am 23.2.2026).
%
%     \bibitem{ruff_linter}
%     ASTRAL, Ruff Linter. In: \url{https://docs.astral.sh/ruff/linter/} (am 23.2.2026).
%
%     \bibitem{alembic_monorepo}
%     FADI BCK, Managing Database Migrations for Multiple Services in a Monorepo with Alembic. In: \url{https://dev.to/fadi-bck/managing-database-migrations-for-multiple-services-in-a-monorepo-with-alembic-3p5l} (am 23.2.2026).
%
%     \bibitem{cron_guru}
%     CRONTAB GURU, Cron schedule expression editor. In: \url{https://crontab.guru/} (am 23.2.2026).
%
%     \bibitem{postgres_pgdump}
%     THE POSTGRESQL GLOBAL DEVELOPMENT GROUP, pg\_dump. In: \url{https://www.postgresql.org/docs/current/app-pgdump.html} (am 23.2.2026).
%
%     \bibitem{postgres_pgrestore}
%     THE POSTGRESQL GLOBAL DEVELOPMENT GROUP, pg\_restore. In: \url{https://www.postgresql.org/docs/current/app-pgrestore.html} (am 23.2.2026).
%
%     \bibitem{mkdocs}
%     MKDOCS, Project documentation with Markdown. In: \url{https://www.mkdocs.org/} (am 23.2.2026).
%
%     \bibitem{mkdocs_material}
%     DONATH Martin, Material for MkDocs. In: \url{https://squidfunk.github.io/mkdocs-material/} (am 23.2.2026).
%
%     % --- Kapitel 7: Qualitätssicherung ---
%     \bibitem{pytest}
%     PYTEST, pytest -- helps you write better programs. In: \url{https://docs.pytest.org/en/stable/} (am 23.2.2026).
%
%     \bibitem{pytest_cov}
%     PYTEST-COV, pytest-cov -- Coverage plugin for pytest. In: \url{https://pypi.org/project/pytest-cov/} (am 23.2.2026).
%
%     \bibitem{postgresql_enum}
%     THE POSTGRESQL GLOBAL DEVELOPMENT GROUP, Enumerated Types. In: \url{https://www.postgresql.org/docs/current/datatype-enum.html} (am 23.2.2026).
%
%     \bibitem{postgresql_datetime}
%     THE POSTGRESQL GLOBAL DEVELOPMENT GROUP, Date/Time Types. In: \url{https://www.postgresql.org/docs/current/datatype-datetime.html} (am 23.2.2026).
%
%     \bibitem{sqlite_datatypes}
%     SQLITE, Datatypes In SQLite. In: \url{https://www.sqlite.org/datatype3.html} (am 23.2.2026).
%
% \end{thebibliography}